\documentclass[10pt,a4paper]{amsart}
\usepackage[margin=19mm]{geometry}
\usepackage{amsmath,amssymb,mathtools}
\usepackage[scr=rsfs,cal=euler]{mathalfa}
\usepackage{xfrac}
\usepackage[unicode,hidelinks,bookmarks=false]{hyperref}
\newcommand{\ie}{i.e.\ }
\newcommand{\cf}{cf.\ }
\newcommand{\resp}{respectively\ }
\newcommand{\Subsection}{\subsection}
\providecommand{\nobreakdash}{\nobreak\hbox{-}}
\setlength{\emergencystretch}{2em}
\multlinegap0pt
\usepackage{amssymb}
\usepackage{enumerate}
\usepackage{mathtools}
\newtheorem{thm}{Theorem}
\newtheorem{prop}{Proposition}
\def\B{\mathbb{B}}
\def\C{\mathbb C}
\def\R{\mathbb R}
\title{Examples of critically cyclic functions in the Dirichlet spaces of the ball}
\author{Pouriya Torkinejad Ziarati}
\date{Source: arXiv:2601.08651v2 (CC BY 4.0)}
\begin{document}
\maketitle

\noindent\emph{Source and licence.} Examples of critically cyclic functions in the Dirichlet spaces of the ball. Version arXiv:2601.08651v2 (CC BY 4.0), distributed by arXiv under the Creative Commons Attribution 4.0 International licence, http://creativecommons.org/licenses/by/4.0/, as recorded in the arXiv licence metadata for this exact version and shown on the abstract page https://arxiv.org/abs/2601.08651v2. Full licence notice: \texttt{licenses/arxiv-2601.08651-CC-BY-4.0.txt}.\par\smallskip

\noindent\textit{Editorial scope.} Counted unit: the article's thm MainThm2 (source0.tex lines 189--198). The article's complete original proof of it is delivered with it (source0.tex lines 666--720, one \texttt{proof} environment of 55 lines, which the article places away from the statement). No mathematics is written, repaired or re-derived: the statement and the proof are the article's own lines, in the article's own notation, and no retained line is substituted or re-flowed.  Retained uncounted as the source-local context the proof needs: lines 161--167; lines 171--180; lines 493--495.  Nothing was omitted from any retained span, so the package carries no omission record.  No retained line contains a citation, so the wrapper carries no bibliography and no bibliography entry.  No retained line contains a figure environment, an image-inclusion command or drawing code, so no image is delivered and the package declares no asset dependency.  Editorial formatting only, in addition: an AMS \texttt{amsart} preamble that restates the article's own declarations and macros used by the retained lines, namely the environments \texttt{thm}, \texttt{prop} on separate counters, together with the standard packages those lines need.  The retained environments print in this document as Theorem 1, Proposition 1 in order of appearance, whereas the article prints the same items with the numbering of its own full document.  The complete native source of the article as distributed in the arXiv e-print for this exact version is bundled unchanged as \texttt{sources/192612/source0.tex}.
        \begin{thm} \label{MainThm1}
Let $E$ be a $\mathcal{K}$-set contained in a simple transverse curve $\Gamma$, and $d \in (0,1)$ be the Hausdorff dimension of $E$. Then $\dim_{K}(E) = d$.

Assume that whenever $t \in \R^+$ is small enough we have $|E_t|  \asymp t^{1 - d}$, and set $\alpha_c = 2 - \dim_{K}(E)$.
Then if  $f$ is an extremal function for the set $E$,
$f$ is cyclic in $D_\alpha(\B_2)$ if and only if $\alpha \leq \alpha_c$.
        \end{thm}

        \begin{prop} \label{MainTool2}
        Let $M$ be a compact $\mathcal{C}^\infty$-smooth complex tangential curve and $E \subseteq M $ be compact. Given $ 0 <\epsilon < \frac{1}{2}$, there exists $f \in \mathcal{O}(\B_2)$ such that $\mathcal{Z}(f) \cap \overline{\B}_2 = E$ and for all $ k$ being a positive integer, we have:
        \begin{equation} \label{eq: Maintool2 estimate1}
             |f(z)| \asymp d_K(z,E)^{1+\epsilon},
        \end{equation}
        and
        \begin{equation} \label{eq: Maintool2 estimate2}
            |R^{k} f(z)| \lesssim d_K(z,E)^{-k+1+\epsilon},
        \end{equation}
        \end{prop}

    \begin{equation} \label{eq:Rg/g_r estimate3}
        \left\| R \frac{g}{g_r} \right\|_{\alpha_c - 2}^2 \lesssim (1-r)^2 \int_{\B_2} \left| \frac{1}{f_r^4(z)}  \right| (1-|z|^2)^{1-\alpha_c}\, \mathrm{d}v(z).
    \end{equation}

        \begin{thm} \label{MainThm2}
        Let $M$ be a compact $\mathcal{C}^\infty$-smooth complex tangential curve, $E \subseteq M $ be compact,
and $d \in (0,1)$ be the Hausdorff dimension of $E$. 
Then $\dim_{K}(E) = \frac{d}{2}$.

Assume further that $|E_t|\asymp t^{1 - d}$; set
$\alpha_c = 2 - \dim_{K}(E)$, and let $f$ be a function
chosen as in Proposition~{\ref{MainTool2}}.
Then $f$ is cyclic in $D_\alpha(\B_2)$ if and only if $\alpha \leq \alpha_c$.
        \end{thm}

\begin{proof} [Proof of Theorem~\ref{MainThm2}]
    Let $f$ be as in Proposition~{\ref{MainTool2}}, using \eqref{eq: Maintool2 estimate1}, \eqref{eq: Maintool2 estimate2} and calculations along the lines of those which led to \eqref{eq:Rg/g_r estimate3} one may show that

    \begin{equation} \label{eq: Rg/g_r estimate2}
    \begin{aligned}
        \left| R \frac{f}{f_r} \right|(z) &\leq  \left|  \frac{(f-f_r) Rf}{f^2_r} \right|(z) + \left|  \frac{ f (Rf - Rf_r) }{f^2_r} \right|(z) \\
        &\lesssim \left|  \frac{(1-r) (Rf)^2}{f^2_r} \right|(z) + \left|  \frac{ (1-r) f R^2f }{f^2_r} \right|(z) \\
        &\lesssim \left|\frac{1-r}{f^2_r(z)}\right| d_K(z,E)^{2 \epsilon}    \\
        &\lesssim \frac{1-r}{d_K(z,E)^2},
    \end{aligned}
    \end{equation}
    and
    \begin{equation} \label{eq: shabih2}
        \left\| R \frac{f}{f_r} \right\|_{\alpha - 2}^2 \lesssim (1-r)^2 \int_{\B_2}  \frac{1}{d_K(rz,E)^{4}}   (1-|z|^2)^{1-\alpha}\, \mathrm{d}v(z).
    \end{equation}
    Also, as before, for $r > \frac{1}{2}$ and $\zeta_0 \in E$ there exists a neighborhood $U$ in $\C^2$ of $\zeta_0$ and a coordinate $Z(x_1,x_2,x_3,x_4)$ valid in $U$ such that for all $z \in \overline{\B}_2 \cap U$ we have
    \begin{equation*}
        d_K(rz,E) \gtrsim \left((1-r)x_1 +|x_2| + x_3^2 + \mathrm{dist}^2(Z(0,0,0,x_4),E)\right).
    \end{equation*}
    Letting $t = \mathrm{dist}(Z(0,0,0,x_4),E)$, $\alpha_c = 2 -\frac{d}{2}$ and noting that we may cover $E$ by finitely many such $U$'s allows us to write
    \begin{equation}
        \begin{aligned}
            \left\| R \frac{f}{f_r} \right\|^2_{\alpha_c - 2}
        &\lesssim (1 - r)^2 \int_{-1}^1 \int_{-1}^1 \int_{-1}^1 \int_0^1 
        \frac{x_1^{1 - \alpha_c}}{((1-r) + x_1 + |x_2| + x_3^2 + t^2)^{4}} \, dx_1\, dx_2\, dx_3\, dx_4 \\
        &\asymp (1 - r)^2 \int_{-1}^1 \int_{-1}^1 \int_{0}^1  
        \frac{1}{((1-r) + x_2 + x_3^2 + t^2)^{2 + \alpha_c}} \, dx_2\, dx_3\, dx_4 \\
        &\asymp (1 - r)^2 \int_{-1}^1 \int_{-1}^1  
        \frac{1}{((1-r) + x_3^2 + t^2)^{1 + \alpha_c}} \, dx_3\, dx_4 \\
        &\asymp (1 - r)^2 \int_{-1}^1  
        \frac{1}{((1-r) + t^2)^{\frac{1}{2} + \alpha_c}} \, dx_4 \\
        &\asymp (1 - r)^2 \int_{-1}^1  
        \frac{t^{2-d}}{((1-r) + t^2)^{\frac{5-d}{2}}} \, dt \asymp 1 < +\infty. \\
        \end{aligned}
    \end{equation}
    Now let $\alpha > \alpha_c$, and $\gamma : [0,1] \rightarrow M$ be a diffeomorphism for $s,s' \in [0,1]$ we have
    \begin{equation*}
        d_K(\gamma(s),\gamma(s')) \asymp |s-s'|^2,
    \end{equation*}
    which implies
    \begin{equation*}
       \dim_{K}(E) = \frac{d}{2}
    \end{equation*}
    and
    \begin{equation*}
            \operatorname{cap}_{\alpha,2}(E) \asymp \operatorname{cap}_{2\alpha-3,1}(\exp{\left(2\pi i \gamma^{-1}(E)\right)}).
    \end{equation*}
    Now similar to Theorem~{\ref{MainThm1}} we have
    \begin{equation*}
        \operatorname{cap}_{\alpha,2}({E}) > 0 
        \quad \Longleftrightarrow \quad
        \int_{0}^{1} \frac{dt}{t^{4 -2\alpha} |E_{t}|} < \infty .
    \end{equation*}
    so $\operatorname{cap}_{\alpha,2}(E)> 0$ and we are done.
\end{proof}

\end{document}
